\chapter{Summary and Future Work}
\label{ch:summary}

Briefly restate the problem you addressed and summarise how your work contributes to solving it. This chapter should give a reader who skips to the end a clear picture of what was done and what comes next.

%%%%%%%%%%%%%%%%%%%%%%%%%%%%%%%%%%%%%%%%%%%%%%%%%%%%%%%%%%%%%%%%%%%%%%%%
\section{Summary}
\label{sec:summary}

Revisit your research questions from Chapter~\ref{ch:introduction} and state concisely how each was addressed. Highlight your key contributions.\footnote{Avoid simply repeating the abstract — this section should reflect on the work as a whole, now that the reader has seen all the details.}

\begin{itemize}
  \item[\textit{RQ1}] How your first research question was answered.
  \item[\textit{RQ2}] How your second research question was answered.
  \item[\textit{RQ3}] How your third research question was answered.
\end{itemize}

%%%%%%%%%%%%%%%%%%%%%%%%%%%%%%%%%%%%%%%%%%%%%%%%%%%%%%%%%%%%%%%%%%%%%%%%
\section{Limitations}
\label{sec:limitations}

Discuss the boundaries of your work — what assumptions were made, what was out of scope, and what threats to validity exist.\footnote{Acknowledging limitations honestly strengthens rather than weakens your thesis.}

%%%%%%%%%%%%%%%%%%%%%%%%%%%%%%%%%%%%%%%%%%%%%%%%%%%%%%%%%%%%%%%%%%%%%%%%
\section{Future Work}
\label{sec:future-work}

Identify open questions and promising directions that follow from your work. What would you do differently with more time? What problems does your work open up that did not exist before?
